
\noindent
\textbf{Egregi Professori Vardanega e Cardin,}

\vspace{0.3cm}

\noindent
Con la presente il gruppo \textbf{Git Happens} intende formalizzare la propria candidatura per la realizzazione del progetto relativo al capitolato \textbf{C0 -- AIRA}, proposto dall'azienda \textbf{Sanmarco Informatica S.p.A.}

\vspace{0.4cm}

\section*{Documentazione di Riferimento}
La valutazione che ha portato a tale decisione si fonda sull'analisi dettagliata e sulla redazione dei seguenti documenti interni ed esterni, consultabili sul \href{https://githappens-team.github.io/#cat-verbali}{sito web ufficiale} del gruppo:

\begin{enumerate}
    \item \textbf{Verbali Interni:}
    \begin{itemize}
        \item Verbale Interno 2026/10/16
        \item Verbale Interno 2026/10/20
        \item Verbale Interno 2026/10/30
    \end{itemize}
    \item \textbf{Verbali Esterni:}
    \begin{itemize}
        \item Verbale Esterno 2026/10/20
        \item Verbale Esterno 2026/10/22
        \item Verbale Esterno 2026/10/26
        \item Verbale Esterno 2026/10/27
    \end{itemize}
    \item \textbf{Analisi dei Capitolati e Preventivo Costi}
\end{enumerate}

\noindent
Per maggiori informazioni relative al confronto tra le diverse proposte e alle motivazioni di dettaglio, si rimanda specificamente al documento \textit{Valutazione Capitolati} e al \textit{Verbale Esterno 2026/10/20}.

\vspace{0.4cm}

\section*{Motivazioni della Scelta}
A seguito dell'analisi preliminare di tutti i capitolati presentati, il gruppo ha individuato in AIRA la proposta più stimolante e maggiormente in linea con i propri interessi formativi e professionali.

Le motivazioni principali che hanno guidato la nostra decisione si articolano sui seguenti punti:
\begin{itemize}
    \item \textbf{Chiarezza e modularità del flusso di lavoro:} L'intero processo risulta lineare e ben definito, dalla gestione dell'acquisizione delle richieste multicanale, all'elaborazione automatizzata mediante Large Language Models (LLM), fino all'esposizione di API REST. La natura fortemente modulare del sistema garantisce a tutti i membri del gruppo l'opportunità di cimentarsi su diversi componenti dello sviluppo software, dall'elaborazione backend dei dati alla realizzazione dell'interfaccia amministrativa.
    
    \item \textbf{Stack tecnologico moderno e solido:} La proposta combina strumenti consolidati nel settore enterprise (\textit{Java} con \textit{Spring Boot} e \textit{Angular}) con tecnologie recenti nel campo dell'Intelligenza Artificiale Generativa (\textit{LangChain4j}, \textit{Whisper} per la trascrizione audio e modelli \textit{OCR} per la lettura dei documenti). Riteniamo che lo studio e la sperimentazione di questi strumenti rappresentino un'importante occasione di crescita. Inoltre, la gestione della sicurezza e l'isolamento multi-tenant dei dati costituiscono una sfida ingegneristica di notevole interesse.
    
    \item \textbf{Disponibilità dell'azienda proponente:} Un elemento determinante è stato il confronto diretto con i referenti di Sanmarco Informatica, i quali si sono mostrati immediatamente disponibili nel chiarire i nostri dubbi iniziali, oltre che estremamente ben organizzati per quanto riguarda il supporto continuo e i futuri incontri di allineamento.
\end{itemize}

\vspace{0.4cm}

\section*{Costi e Tempistiche}
Come dettagliato nel documento \textit{Preventivo Costi}, il gruppo ha stimato l'impegno economico e la data di consegna finale come segue:

\begin{itemize}
    \item \textbf{Scadenza prevista per la consegna:} 19 Marzo 2027
    \item \textbf{Preventivo costi complessivo:} \geneuro\,12.845,00
\end{itemize}

\vspace{0.6cm}

\section*{Composizione del Gruppo}
I componenti del gruppo \textbf{Git Happens} che partecipano alla candidatura:

\vspace{0.3cm}

\noindent
\begin{tabularx}{\textwidth}{@{} X l @{}}
    \toprule
    \textbf{Nome e Cognome} & \textbf{Matricola} \\
    \midrule
    Diana Grosselle & 2137985 \\
    \addlinespace[0.15cm]
    Cristiano Antoniazzo & 2138489 \\
    \addlinespace[0.15cm]
    Gabriele Gambarro & 2147979 \\
    \addlinespace[0.15cm]
    Gianmaria Grinovero & 2138007 \\
    \addlinespace[0.15cm]
    Tommaso Caon & 2111933 \\
    \addlinespace[0.15cm]
    Antonio Hai Xu Tang & 2111017 \\
    \addlinespace[0.15cm]
    Filippo Lissandrin & 2137980 \\
    \bottomrule
\end{tabularx}
\vspace{0.8cm}

\noindent
Restiamo a completa disposizione per qualsiasi chiarimento o integrazione documentale.

\vspace{0.5cm}

\noindent
\textit{Padova, 30 Ottobre 2026}
